\section{From additivity to geodesic currents}

We want to show that if $\nu$ satisfies smoothing, additivity, stability and homogeneity then it comes from an intersection number.
We will follow this strategy:

\begin{enumerate}
    \item Pick a simple closed curve $[x]$, $x \in \pi_1(S,*)$. with base point $p \in [x]$.
    \item For any $c \in \pi_1(S,*)$, crossing $c$ right-handedly (define), we prove
    \[
      \lim_{n \to \infty}  \nu([cx^n])-f([x^n])
    \]
    exists.
    \item For $c,d \in \pi_1(S,*)$, crossing $x$ right-handedly
    
    \[
       \lim_{m,n}  \nu([cx^ndx^m])-\nu([cx^n][dx^m])=0
    \]
    
    \item For $p,q,r,s \in \partial_{\infty} \tilde{S}$, and $a,b,c,d$ translations relating lifts, we set
    
    \[
       \mu([p,q) \times [r,s)) \coloneqq \lim_{n \to \infty} \nu([ax^n][bx^n])-\nu([cx^n][dx^n])
    \]
    
    This exists and moreover is positive, by the condition that $[a] \cup [b]$ smooth to $[c] \cup [d]$.
    \item We define suitable rectangles that form a semi-ring of sets
    \item We prove that the pre-measure we have defined is finitely additive and countably subadditive, and then Caratheodory theorem gives us a measure.
    \item We show that this measure recovers intersection number.
\end{enumerate}

The idea is to reconstruct a current, i.e. a cross-ratio from the periods. (talk/cite previous work by Hamenstadt, Ledrappier, Martone-Zhang?).\\
Work of Otal \cite{Otal90:SpectreMarqueNegative} shows that a current is uniquely determined by its values on closed curves, by using boxes whose endpoints are orbits of certain geodesics.
In his argument, Otal uses lifts of a dense geodesic in the surface, but controlling the asymptotics of those is complicated.
We will instead use boxes whose endpoints are orbits of a closed geodesic $x$, easier to control. The caveat here is that
now we have to be more careful since the target geodesic current may have
delta-masses at the boundaries of the boxes. We will have to specify exactly what's included in the boundary and what is not.
This is where the definition of \emph{right handed} rectangle will come into play. \\

Then we will construct closed curves whose lengths
should correspond to these rectangles. Thus, in addition to the closed geodesic
$x$, we have 4 elements of $\pi_1$, $a,b,c,d$ corresponding to the 4 sides of
these boxes ($c$ and $d$ are opposite sides, and $a,b$ are diagonals), and we will consider limits of

  $$v(a x^n b x^m) - v(c x^n d x^m).$$

One key point here is that the first curve has a crossing that the second one doesn't have
(so the difference is positive). There will be several things to check, all
up to small additive error that goes to zero as $n,m \to \infty$:

\begin{enumerate}
\item $$v(a x^n b x^m) \sim  v(a x^n) + v(b x^m)$$ (same for $c,d$).

\item Both quantities $v(a x^n b x^m)$ and $v(a x^n) + v(b x^m)$ are convex as a function of $n,m$

\item Both quantities depend only on $n+m$
\end{enumerate}


 We use then the additivity of the function $v$ to show it's additive on these boxes, and prove that the boxes form a semi-ring (see proposition \ref{prop:RHsemiring}).
 Then we apply Caratheodory's extension theorem for pre-measures on semi-rings (see theorem \ref{thm:caratheodorymeasures}) to obtain a $\pi_1(S)$-invariant measure.
 Finally, we show these approximations converge to the right value, i.e. to the value of the function on closed curves i.e. the current we have
constructed is really the Liouville current for $v$ (see proposition \ref{prop:liouvillecurrent}).

\subsection{Boxes}
\label{sec:boxes}

Let $v\colon \mathbb{R}\mathcal{C}_+(\Sigma) \to \mathbb{R}_{0 \geq}$ satisfy additive union, smoothing, stability and homogeneity.

We will fix (arbitrarily) a simple closed oriented curve~$x$ that will
serve as endpoints of the allowable rectangles (note that these endpoints are dense in $S^1_{\infty}$, by minimality of the action of $\pi_1(S,*)$ by Deck transformations. See Section~\ref{subsec:actionsfundamental}). For convenience, we will work in $\pi_1(S,\ast)$ with a basepoint~$\ast$ on~$x$, and also
use $x$ for the corresponding element of~$\pi_1$.

We say that $a \in \pi_1(S, \ast)$ \emph{crosses $x$ to the
  right} if there is an essential right-handed crossing.
  This ensures that
  \begin{align}
    \label{eq:cross-red1}[a][x] &\reducesto [ax]
  \end{align}

For $a,b,c,d \in \pi_1(S,\ast)$ curves that cross $x$ to
the right with $a b^{-1} c d^{-1} = 1$, 
let the box $B_{a,b,c,d,x}$
be the \emph{right handed rectangles of type 1 and type 2} defined by the configurations of oriented lifts of $x$ shown in Figures~\ref{fig:box1} and \ref{fig:box2}.

\begin{figure}
\centering{
\resizebox{60mm}{!}{\svgfig{box1}}
\caption{Right handed box Type 1}
\label{fig:box1}
}
\end{figure}

\begin{figure}
\centering{
\resizebox{60mm}{!}{\svgfig{box2}}
\caption{Right-handed box Type 2}
\label{fig:box2}
}
\end{figure}

The elements $a,b,c,d$ in Figure~\ref{fig:box1} are the elements in $\pi_1(S,*)$ relating the lifts of $x$. 
This determines a half-open box $B_{x,a,b,c,d}=[p,q) \times [r,s)$ as illustrated in Figure~\ref{fig:limitingbox}.

\begin{figure}
\centering{
\resizebox{100mm}{!}{\svgfig{limitingbox}}
\caption{End-points of box $B_{x,a,b,c,d}$.}
\label{fig:limitingbox}
}
\end{figure}

\begin{remark}
The box in Figure~\ref{fig:box3} is not a right-handed box. We note that for that box the diagonal component contributes extra because curves transversal to the diagonals are also transverse to one of the lifts of $x$.
\begin{figure}
\centering{
\resizebox{60mm}{!}{\svgfig{box3}}
\caption{Non right-handed box.}
\label{fig:box3}
}
\end{figure}
\end{remark}

We need the curve $[bc]$ to smooth to $[a] \cup [d]$.

Slight issue: Need a little more data to get a box in $\wt S$, rather
than a box up to translation by $\pi_1(S)$. Do this by picking
paths $p_1$, $p_2$, $q_1$, $q_2$ and setting
\begin{align*}
  a &= q_1p_1 & b &= q_2p_1\\
  c &= q_1p_2 & d &= q_2p_2.
\end{align*}

\begin{figure}
\centering{
\resizebox{120mm}{!}{\svgfig{fixedbox}}
\caption{Example of fixed box}
\label{fig:fixedbox}
}
\end{figure}

Explain about whiskers, pointing down or to the right.


\begin{figure}
\centering{
\resizebox{120mm}{!}{\svgfig{whiskers}}
\caption{Example of whiskers}
\label{fig:whiskers}
}
\end{figure}

Allowable boxes
have all corners at whiskers.
\DMGtodo{Need to recover material about whiskers, I forgot many details.}


\subsection{Measures of boxes}

We wish to define a pre-measure by setting
\todo{Labels have been permuted}
\[
\mu_{v}(B_{a,b,c,d,x})=\lim_{n \to \infty}v(bx^n) + v(cx^n)-v(ax^n) - v(dx^m)
\]
To check that this is well-defined, we first prove that this limit
exists.

\begin{lemma}
  If $a$ crosses $x$ to the right, then we have the following
  geometric reductions:
  \begin{align}
    \label{eq:cross-red1}[a x^n][x^m] &\reducesto [a x^{n+m}]\\
    \label{eq:cross-red2}[a x^{n+1}][a x^{n-1}] &\reducesto 2[a x^n]\\
    \label{eq:cross-red3}[a x^n][a x^m] & \reducesto x^{\abs{n-m}}.
  \end{align}
\end{lemma}
\begin{proof}
  Equation~(\ref{eq:cross-red1}) follows since $a$ (and $a x^n)$
  cross~$x$. The other two follow from the bands model.
\end{proof}

\begin{corollary}\label{cor:additivity-limit}
  If $a$ crosses $x$ to the right, then
  \[
    \lim_{n \to \infty} \bigl(v(a x^n) - v(x^n)\bigr)
  \]
  exists.
\end{corollary}

\begin{proof}
  By Equation~(\ref{eq:cross-red2}), $v(a x^n)$ is a convex function
  of $n$, and thus so is $v(a x^n) - v(x^n)$. By
  Equations~(\ref{eq:cross-red1}) and~(\ref{eq:cross-red3}) (both with
  $n=0$),
  \[
    -v(a) \le v(a x^n) - v(x^n) \le v(a).
  \]
  A convex function that is bounded as $n \to \infty$ has a
  well-defined limit.
\end{proof}

Thus the limit defining $\mu_{v}(B_{a,b,c,d,x})$ exists. 